% Part: second-order-logic
% Chapter: syntax-and-semantics
% Section: satisfaction

\documentclass[../../../include/open-logic-section]{subfiles}

\begin{document}

\olfileid[fr]{sol}{syn}{sat}

\olsection{Satisfaction}

\begin{explain}
Pour définir la relation de satisfaction
$\Sat{M}{!A}[s]$ des !!{formula}s du second ordre,
nous devons étendre les définitions aux
!!{variable}s du second ordre. La notion de
!!{structure} reste la même qu'au premier ordre.
Seules les assignations de variables~$s$ changent :
elles doivent désormais donner des valeurs non
seulement aux !!{variable}s du premier ordre,
mais aussi à celles du second ordre.
\end{explain}

\begin{defn}[Assignation de variables]
Une \emph{assignation de variables}~$s$ pour une
!!{structure}~$\Struct{M}$ est une fonction qui
associe à chaque
\begin{enumerate}
\item variable d'objet~$\Obj{v_i}$ un élément
  de~$\Domain{M}$, autrement dit
  $s(\Obj{v_i})\in\Domain{M}$;
\item variable de relation $n$-aire~$\Obj{V_i^n}$
  une relation $n$-aire sur~$\Domain{M}$,
  autrement dit
  $s(\Obj{V_i^n})\subseteq\Domain{M}^n$;
\item variable de fonction $n$-aire~$\Obj{u_i^n}$
  une fonction de $\Domain{M}^n$ dans~$\Domain{M}$,
  autrement dit
  $s(\Obj{u_i^n})\colon\Domain{M}^n\to\Domain{M}$.
\end{enumerate}
\end{defn}

\begin{explain}
Une !!{structure} attribue une !!{value} à chaque
!!{constant} et à chaque !!{function};
une assignation du second ordre attribue des
objets et des fonctions aux variables d'objet et
de fonction. Ensemble, elles permettent de donner
une !!{value} à chaque terme.
\end{explain}

\begin{defn}[Valeur d'un terme]
Si $t$ est un terme du langage~$\Lang L$,
si $\Struct M$ est une !!{structure} pour~$\Lang L$
et si $s$ est une assignation de !!{variable}s
pour~$\Struct M$, sa \emph{!!{value}}
$\Value{t}{M}[s]$ se définit par les clauses du premier ordre, appliquées aux termes du second ordre,
et par la clause supplémentaire suivante :
\begin{quote}
\indcase{t}{\Atom{u}{t_1,\ldots,t_n}}{
\[
\Value{\indfrm}{M}[s]
=s(u)(\Value{t_1}{M}[s],\ldots,
\Value{t_n}{M}[s]).
\]}
\end{quote}
\end{defn}

\begin{defn}[Variante en~$x$]
Si $s$ est une assignation de !!{variable}s pour
une !!{structure}~$\Struct M$, toute assignation
$s'$ pour~$\Struct M$ qui ne diffère de~$s$ que,
éventuellement, par la !!{value} attribuée à~$x$
est une \emph{variante en~$x$} de~$s$.
Nous écrivons alors $\varAssign{s'}{s}{x}$.
On définit de même les variantes en une variable
de relation~$X$ ou de fonction~$u$ du second ordre.
\end{defn}

\begin{defn}
Si $s$ est une assignation de !!{variable}s pour
une !!{structure}~$\Struct M$ et si
$m\in\Domain{M}$, l'assignation
$\Subst{s}{m}{x}$ est définie par
\[
\Subst{s}{m}{x}(y)=\begin{cases}
  m & \text{si }y\ident x,\\
  s(y) & \text{sinon}.
\end{cases}
\]
Si $X$ est une variable de relation $n$-aire et
si $R\subseteq\Domain{M}^n$, l'assignation
$\Subst{s}{R}{X}$ est définie par
\[
\Subst{s}{R}{X}(y)=\begin{cases}
  R & \text{si }y\ident X,\\
  s(y) & \text{sinon}.
\end{cases}
\]
Si $u$ est une variable de fonction $n$-aire et
si $f\colon\Domain{M}^n\to\Domain{M}$,
l'assignation $\Subst{s}{f}{u}$ est définie par
\[
\Subst{s}{f}{u}(y)=\begin{cases}
  f & \text{si }y\ident u,\\
  s(y) & \text{sinon}.
\end{cases}
\]
Dans chacun de ces cas, $y$ peut être n'importe
quelle !!{variable} du premier ou du second ordre.
\end{defn}

\begin{defn}[Satisfaction]
Pour les !!{formula}s du second ordre~$!A$,
la définition de la satisfaction est celle de la
\olref[fol][syn][sat]{defn:satisfaction}, appliquée aux termes et !!{formula}s du second ordre,
augmentée des clauses suivantes :
\begin{enumerate}
\item \indcase{!A}{\Atom{X^n}{t_1,\dots,t_n}}{
  $\Sat{M}{\indfrm}[s]$ si et seulement si
  $\langle\Value{t_1}{M}[s],\dots,
  \Value{t_n}{M}[s]\rangle\in s(X^n)$.}
\tagitem{prvAll}{%
  \indcase{!A}{\lforall[X][!B]}{
  $\Sat{M}{\indfrm}[s]$ si et seulement si, pour
  toute relation $R\subseteq\Domain{M}^n$,
  $\Sat{M}{!B}[\Subst{s}{R}{X}]$.}}{}

\tagitem{prvEx}{%
  \indcase{!A}{\lexists[X][!B]}{
  $\Sat{M}{\indfrm}[s]$ si et seulement s'il existe
  une relation $R\subseteq\Domain{M}^n$ telle que
  $\Sat{M}{!B}[\Subst{s}{R}{X}]$.}}{}

\tagitem{prvAll}{%
  \indcase{!A}{\lforall[u][!B]}{
  $\Sat{M}{\indfrm}[s]$ si et seulement si, pour
  toute fonction $f\colon\Domain{M}^n\to\Domain{M}$,
  $\Sat{M}{!B}[\Subst{s}{f}{u}]$.}}{}

\tagitem{prvEx}{%
  \indcase{!A}{\lexists[u][!B]}{
  $\Sat{M}{\indfrm}[s]$ si et seulement s'il existe
  une fonction $f\colon\Domain{M}^n\to\Domain{M}$
  telle que $\Sat{M}{!B}[\Subst{s}{f}{u}]$.}}{}
\end{enumerate}
\end{defn}

\begin{ex}
Considérons la !!{formula}
$\lforall[z][(\Atom{X}{z}\liff
\lnot\Atom{Y}{z})]$. Elle ne contient aucun
quantificateur du second ordre, mais elle contient
les !!{variable}s du second ordre $X$ et~$Y$,
ici unaires. L'!!{sentence} correspondant du
premier ordre,
$\lforall[z][(\Atom{P}{z}\liff
\lnot\Atom{R}{z})]$, dit que les objets relevant
de l'interprétation de~$P$ ne relèvent pas de celle
de~$R$, et réciproquement. Dans une
!!{structure}, l'interprétation d'un
!!{predicate}~$P$ est donnée par
$\Assign{P}{M}$. En revanche, l'interprétation
des !!{variable}s du second ordre comme $X$
et~$Y$ vient de l'assignation, non de la
!!{structure} elle-même. La !!{formula} du
second ordre n'étant pas un !!{sentence}
(elle contient les !!{variable}s libres $X$
et~$Y$), elle ne peut être satisfaite que
relativement à la fois à une
!!{structure}~$\Struct{M}$ et à une
assignation de !!{variable}s~$s$.

On a
$\Sat{M}{\lforall[z][(\Atom{X}{z}\liff
\lnot\Atom{Y}{z})]}[s]$
si et seulement si chaque !!{element} de~$s(X)$
est hors de~$s(Y)$ et chaque !!{element} hors
de~$s(X)$ appartient à~$s(Y)$; autrement dit,
$s(Y)=\Domain{M}\setminus s(X)$.
Prenons, par exemple,
$\Domain{M}=\{1,2,3\}$. Aucun !!{predicate},
!!{function} ou !!{constant}
n'intervenant ici, seul le domaine de~$\Struct{M}$
importe. Avec $s_1(X)=\{1,2\}$ et
$s_1(Y)=\{3\}$, nous avons
$\Sat{M}{\lforall[z][(\Atom{X}{z}\liff
\lnot\Atom{Y}{z})]}[s_1]$.

En revanche, pour $s_2(X)=\{1,2\}$ et
$s_2(Y)=\{2,3\}$, on a
$\Sat/{M}{\lforall[z][(\Atom{X}{z}\liff
\lnot\Atom{Y}{z})]}[s_2]$.
En effet,
$\Sat{M}{\Atom{X}{z}}[\Subst{s_2}{2}{z}]$,
puisque $2\in\Subst{s_2}{2}{z}(X)$, tandis que
$\Sat/{M}{\lnot\Atom{Y}{z}}
[\Subst{s_2}{2}{z}]$,
puisque $2\in\Subst{s_2}{2}{z}(Y)$ également.
\end{ex}

\begin{ex}
On a
$\Sat{M}{\lexists[Y][(\lexists[y][\Atom{Y}{y}]
\land\lforall[z][(\Atom{X}{z}\liff
\lnot\Atom{Y}{z})])]}[s]$
s'il existe un ensemble $N\subseteq\Domain{M}$
tel que
$\Sat{M}{(\lexists[y][\Atom{Y}{y}]
\land\lforall[z][(\Atom{X}{z}\liff
\lnot\Atom{Y}{z})])}[\Subst{s}{N}{Y}]$.
Cela se produit exactement lorsque $N$ est non
vide, de sorte que
$\Sat{M}{\lexists[y][\Atom{Y}{y}]}
[\Subst{s}{N}{Y}]$, et que, comme dans l'exemple
précédent, $N=\Domain{M}\setminus s(X)$.
Autrement dit,
$\Sat{M}{\lexists[Y][(\lexists[y][\Atom{Y}{y}]
\land\lforall[z][(\Atom{X}{z}\liff
\lnot\Atom{Y}{z})])]}[s]$
si et seulement si
$\Domain{M}\setminus s(X)\neq\emptyset$,
ou encore $s(X)\neq\Domain{M}$.
La !!{formula} est donc satisfaite, par exemple,
si $\Domain{M}=\{1,2,3\}$ et
$s(X)=\{1,2\}$, mais pas si
$s(X)=\{1,2,3\}=\Domain{M}$.

Puisque la !!{formula} n'est pas satisfaite lorsque
$s(X)=\Domain{M}$, l'!!{sentence}
\[
\lforall[X][\lexists[Y][(\lexists[y][\Atom{Y}{y}]
\land\lforall[z][(\Atom{X}{z}\liff
\lnot\Atom{Y}{z})])]]
\]
n'est jamais satisfait : dans toute
!!{structure}~$\Struct{M}$, l'instance où~$X$
reçoit $\Domain{M}$ rend la portée du quantificateur
universel fausse. En revanche, l'!!{sentence}
\[
\lexists[X][\lexists[Y][(\lexists[y][\Atom{Y}{y}]
\land\lforall[z][(\Atom{X}{z}\liff
\lnot\Atom{Y}{z})])]]
\]
est satisfait relativement à toute
assignation~$s$, car le domaine, non vide, possède
toujours un sous-ensemble propre~$R$
(par exemple, $R=\emptyset$).
\end{ex}

\begin{ex}
L'!!{sentence} du second ordre
$\lforall[X][\lforall[y][X(y)]]$ affirme que
toute relation unaire, c'est-à-dire toute propriété,
vaut de chaque objet. Cela n'est jamais vrai :
dans toute !!{structure}~$\Struct{M}$ de domaine
non vide, on peut choisir une assignation~$s$
telle que $s(X)=\emptyset$ et
$s(y)=a\in\Domain{M}$; alors
$\Sat/{M}{X(y)}[s]$.
Ainsi, pour toute !!{formula}~$!A$ et toute
assignation~$s$,
$\Sat{M}{!A\lif
\lforall[X][\lforall[y][X(y)]]}[s]$
si et seulement si $\Sat{M}{\lnot !A}[s]$.
Autrement dit, au second ordre, on peut définir
$\lnot$ au moyen de $\lforall$ et de~$\lif$.
\end{ex}

\begin{prob}
Montrez qu'en logique du second ordre,
$\lforall$ et~$\lif$ permettent de définir
les autres connecteurs :
\begin{enumerate}
\item Soit $X$ une variable de relation unaire qui n'est libre ni dans~$!A$ ni dans~$!B$. Démontrez que
  $!A\land !B$ équivaut en logique du second ordre à
  $\lforall[X][((!A\lif
  (!B\lif\lforall[x][X(x)]))\lif
  \lforall[x][X(x)])]$.
\item Trouvez une !!{formula} du second ordre
  n'utilisant que $\lforall$ et $\lif$ et
  équivalente à $!A\lor !B$.
\end{enumerate}
\end{prob}

\begin{editorial}
Dans les trois définitions d'assignations modifiées,
la source écrit respectivement
$\Subst{s}{m}{y}$, $\Subst{s}{M}{y}$ et
$\Subst{s}{f}{y}$ au membre de gauche, alors que
la variable modifiée est~$x$, $X$ ou~$u$ et que
$y$ est l'argument auquel on applique
l'assignation. Nous écrivons donc
$\Subst{s}{m}{x}(y)$,
$\Subst{s}{R}{X}(y)$ et
$\Subst{s}{f}{u}(y)$, en réservant~$R$ à la
relation qui remplace la valeur de~$X$.
Dans la dernière équivalence de satisfaction,
nous indiquons aussi l'assignation~$s$, puisque
$!A$ peut contenir des !!{variable}s libres.
\end{editorial}

\end{document}
